\hypertarget{control-de-flujo}{}
\pdfbookmark[0]{Control de flujo}{control-de-flujo}
\begin{frame}\frametitle{Control de flujo}
\begin{itemize}
\item {} 
No hay diferencias con otros lenguajes de programación
\emph{interpretados}
\item {} 
En estos lenguajes un bucle es en realidad un \emph{iterador}
\end{itemize}

\end{frame}

% ___________________________________________________________________________

\hypertarget{condicionales}{}
\pdfbookmark[0]{Condicionales}{condicionales}
\begin{frame}[containsverbatim]
\frametitle{Condicionales}

Este es un ejemplo del uso de la sentencia \emph{if}:
\begin{Octave}
if saludo 
  disp('hola') 
else 
  disp('no te saludo') 
end
\end{Octave}

\begin{itemize}
\item 
  ¿Cuál es la salida si saludo=1?
\item 
  En Octave, podemos usar \texttt{endif} (no compatible Matlab!).
  Igual ocurrirá con \texttt{endfor}, \texttt{endfunction}
\end{itemize}
\end{frame}

% ___________________________________________________________________________

\hypertarget{bucles}{}
\pdfbookmark[0]{Bucles}{bucles}
\begin{frame}[containsverbatim]
\frametitle{Bucles}

Ejemplo del uso de la sentencia \emph{for}:
\begin{Octave}
function primetest(n) 
  printf('Numeros primos de 1 a %i\n',n) 
  for i=1:n 
    if isprime(i) 
      disp(i) 
    endif
  endfor
endfunction
\end{Octave}
\textbf{Ejercicio}:
\begin{itemize}
\item Programar la función en el fichero
  ``\texttt{primetest.m}''
\item 
  ¿Cuál es la salida de esta subrutina si n=20?
\end{itemize}
\end{frame}

% ___________________________________________________________________________

\hypertarget{bucles-iteradores}{}
\pdfbookmark[0]{Bucles iteradores}{bucles-iteradores}
\begin{frame}[containsverbatim]
\frametitle{Bucles iteradores}

Recordad que secuencia == vector
\begin{Octave}
>> for i=[1,3,2,4,3,4,5] 
> disp(i) 
> end 
 1 
 3 
 2 
 4 
 3 
 4 
 5
\end{Octave}

\end{frame}

% ___________________________________________________________________________

\hypertarget{otras-sentencias}{}
\pdfbookmark[0]{Otras sentencias}{otras-sentencias}
\begin{frame}\frametitle{Otras sentencias}
\begin{quote}
\begin{description}

\item [while:]
Bucle controlado por condición lógica

\item [case:]
Control de casos cerrados

\item [try:]
Sentencia de control para probar errores

\item [break:]
Permite abandonar directamente el interior de un bucle


\item [continue:]
Se utiliza para saltar al final de la actual iteración de un bucle


\item [return:]
Devuelve el control al programa principal

\end{description}
\vfill
\textbf{Ejercicio} Realizar distintos ejemplos con las
sentencias anteriores
\end{quote}

\end{frame}



% ___________________________________________________________________________

\begin{frame}[fragile]{Funciones en MATLAB}


\lstset{
  language=Matlab,
  basicstyle=\ttfamily\footnotesize,
  keywordstyle=\color{blue!70!black}\bfseries,
  commentstyle=\color{green!40!black}\itshape,
  frame=single,
  framerule=0.3pt,
  columns=fullflexible,
  keepspaces=true,
  showstringspaces=false,
  numbers=none
}

\begin{columns}[T]
  \begin{column}{0.52\textwidth}
    \textbf{1. Función en fichero} (\texttt{\small area\_circulo.m})
\begin{lstlisting}
function [A, P] = area_circulo(r)
% AREA_CIRCULO  Area y perimetro
%   [A, P] = area_circulo(r)
  A = pi * r.^2;
  P = 2 * pi * r;
end
\end{lstlisting}
    \vspace{0.4em}
    \textbf{Uso:}\vspace{-0.5em}
\begin{lstlisting}
[A, P] = area_circulo(2);
A = area_circulo(1:3);
\end{lstlisting}
  \end{column}

  \begin{column}{0.46\textwidth}
    \textbf{2. Función anónima}
\begin{lstlisting}
f = @(x) x.^2 + 1;
f(3)       % 10
a = 2;
g = @(x) a*x;  % captura a
\end{lstlisting}
    \vspace{0.4em}
    \textbf{Como argumento:}
\begin{lstlisting}
fzero(@(x) cos(x)-x, 0.5)
integral(@(x) exp(-x.^2), 0, 1)
\end{lstlisting}
  \end{column}
\end{columns}

\vspace{0.6em}
\begin{itemize}
  \item El fichero debe llamarse \texttt{area\_circulo.m}: manda el nombre del fichero.
  \item Varias salidas entre corchetes; \texttt{nargin} y \texttt{nargout} dan el n.º de argumentos.
  \item Usa \texttt{.*}, \texttt{./}, \texttt{.\^{}} para que la función admita vectores.
  \item Una anónima es una sola expresión y congela las variables externas al definirla.
\end{itemize}
\end{frame}


% ___________________________________________________________________________

\hypertarget{funciones}{}
\pdfbookmark[0]{Funciones}{funciones}
\begin{frame}[containsverbatim]
  \frametitle{Las funciones con algo más de profundidad}

Varios valores devueltos
\begin{Octave}
function [m,M] = minmax(valores) 
  % Esta funci'on toma como par'ametro un vector y 
  % devuelve sus valores m'inimo y m'aximo
  m = min(valores); 
  M = max(valores);
endfunction 
 
help minmax
lista = 1./(1:10);
m = minmax(lista);
\end{Octave}

\textbf{Observaciones}
\begin{itemize}
\item Recordar que \texttt{endfunction} no es compatible con Matlab
\item \texttt{endfunction} es opcional si función está en su fichero 
\end{itemize}

\end{frame}

\begin{frame}[fragile]
\frametitle{Ejercicio}
\begin{itemize}
    \item Programa una función llamada \texttt{integracion\_numer(f, a, b, N)}, que aplique la fórmula de los trapecios compuesta\footnote{\url{https://es.wikipedia.org/wiki/Regla_del_trapecio\#Regla_del_trapecio_compuesta}} con \texttt{N} subintervalos para aproximar la integral de una función $f$ en un intervalo $(a,b)$.
  \item Calcula de forma manual la integral de algunas funciones sencillas, por ejemplo $f_1(x)=x^2$, $f_2(x)=\cos(x)$, $f_3(x)=3\cos(3x)\cdot e^{\sin(3x)}$.
  \item Aproxima estas integrales con la función que hemos definido con diferentes valores de $N$ y evalúa el error respecto a la interal exacta que calculaste
  \item Genera una tabla con los resultados
\end{itemize}
\end{frame}

% ___________________________________________________________________________

\hypertarget{funciones-an-nimas}{}
\pdfbookmark[0]{Funciones anónimas.}{funciones-an-nimas}
\begin{frame}[containsverbatim]
\frametitle{Más sobre funciones anónimas}
\small
Definición
\begin{itemize}
\item Matlab como casi todos los senguajes, define un tipo especial de funciones,
  (\emph{funciones lambda}). 
 Para ello, se utiliza el carácter \verb|@|
\item Permite definr funciones breves sin necesidad de un archivo
  adicional\footnote{Matlab no permite la definición de funciones de
    forma interactiva (salvo \texttt{inline}). Octave sí lo permite.}

\end{itemize}
\begin{Octave}
>> testf = @(x,y) exp(-(x.^2+y.^2)) 
>> testf(1,i) 
ans = 1
\end{Octave}
Utilidad: estas funciones  pueden pasarse como argumento para otra
función
\begin{Octave}
>> integracion_numer(@(x) x^2, 0, 1);
\end{Octave}
\end{frame}


% ___________________________________________________________________________

\begin{frame}[containsverbatim]
\frametitle{Problemas arreglados por Octave}

Declaración de funciones directamente en consola:
\begin{Octave}
>> function y=foo(x) 
> y=3*x 
> endfunction
>> foo(4) 
y =  12 
ans =  12
\end{Octave}
\end{frame}

\begin{frame}[containsverbatim]
 \frametitle{Problemas arreglados por Octave}
 Varias funciones en un mismo archivo.
 \begin{itemize}
 \item 
 Definimos el fichero \texttt{biblioteca.m}
 \begin{Octave}
   function y=cuadrado(x)
     y=x^2;
   end
   function y=raiz(x)
     y=sqrt(x);
   end
 \end{Octave}
 \item 
 Podemos acceder a las dos funciones:
 \begin{Octave}
   >> source('biblioteca.m')
   >> cuadrado(raiz(2))
   ans =  2.0000
 \end{Octave}
\end{itemize}
\end{frame}



% ___________________________________________________________________________
\begin{frame}[containsverbatim]
\frametitle{Problemas arreglados por Octave}

Ejecución de ``funciones compuestas''
\begin{Octave}
>> a=rand(3,3)(1,2) 
a =  0.14272
\end{Octave}

Permite expresiones mucho más compactas.

\end{frame}


% ___________________________________________________________________________

\hypertarget{la-biblioteca-de-funciones}{}
\pdfbookmark[0]{La biblioteca de funciones}{la-biblioteca-de-funciones}
\begin{frame}\frametitle{La biblioteca de funciones}
\begin{itemize}
  \setlength{\itemsep}{0.7em}
\item {} 
  Integración numérica (\textbf{quad}, quadl, quad2dg...)

\item {} 
  Interpolación (\textbf{interp1}, interp2, polyfit, fft,...)

\item {} 
  Resolución numérica de ecuaciones diferenciales (\textbf{ode45},...)

\item {} 
Estadística descriptiva (mean, std, median,...)

\end{itemize}

% \textbf{Ejercicio}: Localiza funciones de tu interés y pruébalas!
% \begin{itemize}\item 
%   Indicación:
%   \begin{itemize}
%   \item Echa un vistazo a la ayuda de Octave (HTML)
%   \item Utiliza la función \texttt{lookfor} (busca funciones relacionadas
%     con una palabra clave). Por ejemplo, \texttt{lookfor fourier}.
%   \end{itemize}

% \end{itemize}


\end{frame}

% ___________________________________________________________________________

% \hypertarget{toolkits-en-octave}{}
% \pdfbookmark[0]{``Toolkits'' en Octave}{toolkits-en-octave}
% \begin{frame}\frametitle{``Toolkits'' en Octave}
% \begin{itemize}
% \item {} 
% \texttt{Tookits} ``oficiales'': Octave forge \href{http://octave.sourceforge.net}{http://octave.sourceforge.net}
% \begin{itemize}
% \item Mirar en la web la lista de \texttt{toolkits}
% \end{itemize}
% \item {} 
% Muchos otros.
% \begin{quote}
% \begin{itemize}
%
% \item MPITB (\texttt{Toolkit} de Paralelización con MPI, desarrollado por
%   Javier F. Baldomero, Universidad de Granada)\footnote{\url{http://atc.ugr.es/javier-bin/mpitb}}
%
% \item {} 
% Octaviz (VTK en Octave)
%
% \item {} 
% Octplot, etc
%
% \end{itemize}
% \end{quote}
%
% \item Problema: no hay garantías de compatibilidad con \texttt{Toolkits} de Matlab
%   \begin{itemize}
%   \item Independencia respecto a Octave y sus líneas estratégicas
%   \end{itemize}
% \end{itemize}
%
% \end{frame}
%
% % ___________________________________________________________________________

% \hypertarget{ejercicios}{}
% \pdfbookmark[0]{Ejercicios}{ejercicios}
% \begin{frame}\frametitle{Ejercicios}

% \end{frame}

% % ___________________________________________________________________________

% \hypertarget{ejercicio}{}
% \pdfbookmark[0]{Ejercicio}{ejercicio}
% \begin{frame}\frametitle{Ejercicio}

% Resolver la ecuación de Laplace en 1d con CC Dirichlet:

% Mediante diferencias finitas.

% \end{frame}

% % ___________________________________________________________________________

% \hypertarget{id6}{}
% \pdfbookmark[0]{Ejercicio}{id6}
% \begin{frame}\frametitle{Ejercicio}

% \end{frame}

% % ___________________________________________________________________________

% \hypertarget{id7}{}
% \pdfbookmark[0]{Ejercicio}{id7}
% \begin{frame}\frametitle{Ejercicio}

% Crear una \emph{función} que genere esta matriz con \emph{n} arbitrario

% Usando también la función \emph{diag}

% \end{frame}

% % ___________________________________________________________________________

% \hypertarget{id8}{}
% \pdfbookmark[0]{Ejercicio}{id8}
% \begin{frame}\frametitle{Ejercicio}

% Crear una función que cree el término inEependiente con \emph{n}
% arbitrario:

% y resolver para \emph{n} = 7

% \end{frame}

% ___________________________________________________________________________

% \begin{frame}[containsverbatim]
% \frametitle{Extensiones dinámicas de Octave}
% \vspace{-0.5em}
%   Fichero \texttt{MiFuncion.cpp}:
%
%   \begin{C++}
% #include <oct.h> // Acceso al codigo de Octave!
% DEFUN_DLD(MiFuncion, args, nargout, "Descripcion")  
% {
%   // ... Codigo C++ ...
%   return octave_value()
% }
%   \end{C++}
%
%   Compilación (con \texttt{mkoctfile}, un envoltorio del compilador  C++)
%
%   \begin{sh}
% $ mkoctfile MiFuncion.cpp
%   \end{sh} %$
%
%   Obtenemos el fichero \texttt{MiFuncion.oct}, que usaremos desde octave:
%   \begin{Octave}
% >> MiFuncion
%   \end{Octave}
% \end{frame}
%
%
% \begin{frame}[containsverbatim]
% \frametitle{Extensiones dinámicas de Octave: Hola Mundo!}
%
% Fichero \texttt{holamundo.cpp}
% \begin{C++}
% #include <octave/oct.h>
% DEFUN_DLD(holamundo, args, nargs , ``Programa hola mundo'')
% {
%   std::cout << ``Hola Mundo!'' << std::endl;
%   return octave_value()
% }
% \end{C++} 
%
% Lo compilamos (\verb|$ mkoctfile holamundo.cpp|) y lo usamos desde Octave:
%   \begin{Octave}
% >> holamundo
% ``Hola Mundo!''
%   \end{Octave}
% \end{frame}
%
% \begin{frame}[containsverbatim]
% \frametitle{Extensiones dinámicas de Octave: $f(x)=x^2$ }
% \small
% Fichero \texttt{xcuadrado.cpp}:
% \begin{C++}
% #include <octave/oct.h>
% DEFUN_DLD(xcuadrado, args, , \
%   "Calcula el cuadrado de un real, x'')
% {
%   if(args.length() != 1) {
%     std::cerr << "ERROR: Se necesita tomar un valor" << std::endl;
%     return octave_value()
%   }
%   const double x = args(0).double_value();
%   octave_value y = x*x;
%   return y;
% }
% \end{C++}
% \end{frame}
%
% \begin{frame}[containsverbatim]
%   \frametitle{Creación de módulos dinámicos C/C++ con SWIG} 
%
%   \begin{itemize}
%   \item 
%   \textbf{¿Qué es SWIG?} 
%   Potente interfaz para generar envolturas (\textit{wrappers}) que
%   permiten acceder a funciones (o clases) C/C++ desde Octave, Python,
%   Java, Perl,...
%   \item \textbf{¿Cómo usar SWIG?}
%     \begin{enumerate}
%     \item Tenemos varias funciones (o clases) C/C++ en un fichero
%       ``\verb|ejemplo.cpp|''
%     \item Escribimos un fichero de directivas ``\verb|ejemplo.i|'', donde
%       indicamos a qué funciones C/C++ deseamos acceder
%     \item A partir de él, SWIG genera un fichero C/C++ ``\verb|ejemplo_wrap.cpp|''
%     \item Compilamos ``\verb|ejemplo_wrap.cpp|'' $+$
%       ``\verb|ejemplo.cpp|'' (mediante \verb|mkoctfile|), obteniendo
%       ``\verb|ejemplo.oct|''
%     \item Desde Octave, ya podemos utilizar el módulo o \emph{toolkit}
%       ``\verb|ejemplo|'' 
%     \end{enumerate}
%   \end{itemize}
% \end{frame}
% ___________________________________________________________________________

\begin{frame}\frametitle{Bibliografía}
\begin{itemize}
\item {}
\textit{Introducción informal a Matlab y Octave}:\par
\href{http://iimyo.forja.rediris.es/}{http://iimyo.forja.rediris.es/}

\item {}
\emph{A. Quarteroni, F. Salieri} Cálculo científico con Matlab y
Octave. Springer 2006.

\item {}
\textit{Documentación ``oficial'':}\par
\href{http://www.gnu.org/software/octave/docs.html}{http://www.gnu.org/software/octave/docs.html}

\item {} 
\textit{Documentación de Octave-Forge:}\par
\href{http://octave.sourceforge.net/docs.html}{http://octave.sourceforge.net/docs.html}
\begin{itemize}
\item {} 
Incluye
\href{http://octave.sourceforge.net/coda/coda.pdf}{http://octave.sourceforge.net/coda/coda.pdf},
manual para extensiones dinámicas en C++.

\end{itemize}

\end{itemize}
\end{frame}

%%% Local Variables: 
%%% mode: latex
%%% TeX-master: "modulo_octave"
%%% End: 
